% titel: Blatt 3, A8: Wärmekapazität eines Kalorimeters
% skript: Mortimer Kap. 5.3 (Kalorimetrie, Wärmekapazität)

% Teil I: Werkzeuge, kein Lösungsweg.

Bei der Kalorimetrie misst man die Temperaturänderung, die eine zugeführte
Wärmemenge im Kalorimeter hervorruft. Die Wärmekapazität verbindet beide
Größen.

\begin{werkzeugkasten}[Wärmekapazität des Kalorimeters]
Für ein abgeschlossenes Kalorimeter gilt näherungsweise:
\begin{equation}
  Q = C\cdot \Delta T, \qquad C = \frac{Q}{\Delta T}.
\end{equation}
$Q$ = zugeführte Wärme, $\Delta T = T_{\text{end}} - T_{\text{anf}}$
(in Kelvin oder in °C, die Differenz ist gleich), $C$ = totale
Wärmekapazität. Man kann $C$ schreiben als $C = C_{\text{Gerät}} +
m_{\text{Wasser}}\cdot c_p(\mathrm{H_2O})$; hängt aber in der Aufgabe keine
Wassermasse an, so genügt $C = Q/\Delta T$ direkt.

\textit{Skript:} Mortimer Kap. 5.3 (Kalorimetrie), Übung 3/8.
\end{werkzeugkasten}

\begin{achtung}[Darauf achten]
\begin{itemize}
  \item $\Delta T$ als Differenz $T_{\text{end}} - T_{\text{anf}}$ bilden;
        für die Differenz ist es gleich, ob in °C oder K gerechnet wird.
  \item Energie und Wärmekapazität in passenden Einheiten: $Q$ in J, $C$ in
        J/K; oder beides in kJ (dann J/K $\to$ kJ/K umrechnen).
  \item Wärmekapazität ist proportional zur \emph{gesamten} erwärmten
        Materie; je grösser $C$, desto kleiner die Temperaturänderung bei
        gleicher Wärme.
  \item Mehr Konstanten als nötig sind eine Falle: hier sind
        $c_p(\mathrm{H_2O})$ und $\rho(\mathrm{H_2O})$ erst dann nötig, wenn
        eine Wassermasse vorliegt.
\end{itemize}
\end{achtung}
